%% kapa: exam=1 round=37 year=2026 period=2 form=A subject=appraisal_law subject_ko=감정평가관계법규 seq=37 number=37
%% source: Q-Net 감정평가사 기출문제 (cst00309) artlSeq=5258924 file=raw/1008/5258924/2026년 제37회 감정평가사 1차 2교시.pdf
%% conversion: kordoc 4.24.51 -> tools/gen1.py (auto); review=proofread-sonnet; flags=none
\begin{question}{37}
도시 및 주거환경정비법령상 조합에 관한 설명으로 옳지 않은 것은?
\begin{choices}
  \item 조합임원의 임기는 3년 이하의 범위에서 정관으로 정하되, 연임할 수 있다.
  \item 조합장 또는 이사가 자기를 위하여 조합과 소송을 할 때에는 주민대표회의 위원장이 조합을 대표한다.
  \item 조합임원의 임기만료 후 6개월 이상 조합임원이 선임되지 아니한 경우에는 시장ㆍ군수등이 조합임원 선출을 위한 총회를 소집할 수 있다.
  \item 조합장은 선임일부터 관리처분계획인가를 받을 때까지는 해당 정비구역에서 거주(영업을 하는 자의 경우 영업을 말함)하여야 한다.
  \item 정관으로 정하는 바에 따라 대의원회 소집청구가 있는 때에는 조합장은 해당일부터 14일 이내에 대의원회를 소집하여야 한다.
\end{choices}
\end{question}
\answer{2}
